% APF v36 -- Derivation of the Trace-to-Scheme Export Gates from APF
\section{Derivation of the Trace-to-Scheme Export Gates from APF}
\label{sec:g1-g6-apf-derivation}

\paragraph{Purpose.}
The preceding Trace-to-Scheme Export Theorem states six gates for promoting an APF trace anchor to a physical export candidate.  This section derives those gates from APF itself.  The gates are not external reviewer hygiene.  They are the necessary conditions for a trace-to-observable comparison to be an admissible continuation between contexts.

\subsection{APF source principles used}

We use only the following APF spine principles.

\begin{enumerate}
\item \textbf{Contextual continuation.}  Physical meaning is expressed by judgments of the form
\[
        \Gamma\vdash_C D\to E,
\]
meaning that, in context $\Gamma$, distinction-expression $D$ can continue as $E$ using enforcement capacity no greater than $C$.

\item \textbf{Meaning by continuation profile.}  A distinction-expression has physical meaning only through its admissible continuation profile
\[
        \operatorname{Cont}_{\Gamma}(D)
        =\{E:\exists C<\infty\text{ with }\Gamma\vdash_C D\to E\}.
\]
Two expressions that cannot be separated by any admissible tested continuation are physically equivalent at the tested resolution.

\item \textbf{Finite cost-bearing distinctions.}  A physical distinction is a finite, positive-cost separator of configuration profiles.  Thus a physical comparison must specify what separation is being tested, in what context, and at what finite resolution.

\item \textbf{Specificity of commitments.}  Distinctions are specific structural commitments.  Their cost depends on the partition, scale, convention, and continuation route by which they are carried.

\item \textbf{Closed ledger.}  In a finite physical regime the enforcement ledger closes: no distinction, convention, route assumption, uncertainty choice, or residual may be silently placed off-ledger.

\item \textbf{No target reuse.}  A claimed prediction or validation must be carried by a continuation that is independent of the target distinction being predicted or validated.  If the target observable is consumed as an input to the route, the comparison is a fitted relabeling, not an APF export.
\end{enumerate}

\begin{definition}[Trace context and physical codomain context]
A \emph{trace anchor} $T$ is a locally closed APF distinction-value in an internal trace context $\Gamma_T$.  A \emph{physical observable} $O_S$ is a tested distinction in an external measurement-scheme context $\Gamma_S$, where $S$ names the scheme, convention, scale, and observable type.  Examples include on-shell electroweak, self-scale $\overline{\rm MS}$, pole, MSR, Monte Carlo, and lattice $\overline{\rm MS}(2\,\mathrm{GeV})$ codomains.
\end{definition}

\begin{definition}[APF export]
An APF export is an admissible comparison continuation from the trace context to the physical codomain context:
\[
        \Gamma_{T\to S}\vdash_C T\to O_S.
\]
Equivalently, an export is witnessed by a route tuple
\[
        \mathcal R=(T,S,\Phi,\Lambda,\Sigma,\mathcal A,\mathcal C_{\rm res}),
\]
where $S$ is the codomain, $\Phi$ is the transport map or justified identity admission, $\Lambda$ is the constants/scale/convention ledger, $\Sigma$ is the uncertainty--covariance protocol, $\mathcal A$ is the no-smuggling audit, and $\mathcal C_{\rm res}$ assigns any residual to named route channels.
\end{definition}

\begin{theorem}[Export gates are APF-derived]
\label{thm:gates-derived-from-apf}
If a trace anchor $T$ is promoted to a physical export candidate for observable $O_S$, then gates G1--G6 follow from the APF spine principles above.
\end{theorem}

\begin{proof}
Promotion to a physical export candidate means that the comparison $T\leadsto O_S$ is not merely a numerical proximity claim.  It must be an APF-meaningful continuation from an internal trace context to an external measurement-scheme context:
\[
        \Gamma_{T\to S}\vdash_C T\to O_S.
\]
Each gate is a necessary component of that judgment.

\smallskip
\noindent\textbf{G1: Codomain declared.}
By contextual continuation and meaning by continuation profile, a distinction-expression has physical meaning only in a context.  The symbol $O$ without a scheme $S$ does not determine a continuation profile: the same numerical-looking quantity may denote inequivalent tested distinctions in on-shell, $\overline{\rm MS}$, pole, MSR, Monte Carlo, or lattice contexts.  Therefore the target must be typed as $O_S$ before the judgment $T\to O_S$ is well formed.  Hence codomain declaration is forced.

\smallskip
\noindent\textbf{G2: Transport map supplied.}
The APF primitive is not equality of scalars; it is admissible continuation.  To assert export from $T$ to $O_S$ one must supply a witness that carries the trace distinction into the codomain distinction:
\[
        \Phi:T\mapsto O_S,
\]
or else justify that $T$ already lies in $S$ by identity codomain admission.  Without $\Phi$ there is no continuation witness, only an external comparison of symbols.  The map must be non-inverse with respect to the target observable: if $\Phi$ is constructed by consuming $O_S$, then $T\to O_S$ is not a source-independent continuation.  Hence a non-inverse transport map or identity admission is forced.

\smallskip
\noindent\textbf{G3: Constants ledger.}
By finite cost-bearing distinction and specificity of commitments, the route $\Phi$ is not free.  Its scales, external constants, renormalization conventions, row choices, thresholds, and normalization choices are themselves structural commitments required to carry the comparison.  By closed-ledger accounting, every such commitment must be placed on-ledger:
\[
        \Lambda(\Phi)=\{\text{constants, scales, conventions, thresholds, normalizations}\}.
\]
If these are omitted, information has entered the comparison through an unpriced side channel.  Hence a constants ledger is forced.

\smallskip
\noindent\textbf{G4: Covariance protocol.}
A physical comparison is a finite tested distinction, not an exact real-number identity outside resolution.  The route must say how uncertainty in $T$, in $\Lambda$, and in the transport procedure propagates to uncertainty in $O_S$:
\[
        \Sigma_{O}=D\Phi\,\Sigma_{(T,\Lambda)}\,D\Phi^{\!*}+\Sigma_{\rm route}
\]
where the displayed expression is schematic and may be replaced by the appropriate finite-row, nonlinear, Monte Carlo, interval, or covariance evaluator.  Without such a protocol the residual has no tested meaning: one cannot tell whether the comparison is within the finite resolution of the route or outside it.  Since APF distinctions are finite cost-bearing separators at finite resolution, uncertainty/covariance propagation is forced.

\smallskip
\noindent\textbf{G5: No-smuggling audit.}
By meaning through admissible continuation, an export route must distinguish a source-carried continuation from a target-fitted relabeling.  If the target observable is used to choose $\Phi$, $\Lambda$, $\Sigma$, or a normalization parameter, then the alleged prediction no longer has an independent continuation profile.  The target has been placed inside the route that is supposed to reach it.  This violates closed-ledger accounting and collapses validation into inverse fitting.  Hence a no-smuggling audit is forced.

\smallskip
\noindent\textbf{G6: Residual channels assigned.}
After a declared route is applied, any discrepancy
\[
        R=O_S-\Phi(T)
\]
is itself a distinction: it says what remains uncarried by the route.  By closed-ledger accounting, that distinction cannot be hidden inside the trace anchor or ignored.  It must be assigned to a named channel: transport truncation, scheme mismatch, missing threshold, QCD/chiral/lattice obstruction, covariance, residual generation operator, or new physics.  Otherwise the comparison uses an unaccounted reservoir to absorb mismatch.  Hence residual-channel assignment is forced.

\smallskip
Together G1--G6 are exactly the data required for $T\to O_S$ to be an APF-admissible comparison continuation.  Therefore the gates are consequences of APF's own spine, not additional empirical conventions.  \qedhere
\end{proof}

\begin{corollary}[Gate failure is an APF obstruction, not trace failure]
If any one of G1--G6 fails, the trace anchor may remain closed in its internal APF context $\Gamma_T$, but the export judgment
\[
        \Gamma_{T\to S}\vdash_C T\to O_S
\]
is not established.  The correct status is therefore $[P_{\rm trace}]$, $[P_{\rm local}]$, or $[P_{\rm obstruction}]$, not $[P_{\rm physical\ final}]$.
\end{corollary}

\begin{proof}
Failure of a gate removes one component required to form or witness the comparison continuation.  G1 failure leaves the target untyped.  G2 failure leaves the continuation unwitnessed.  G3 failure places route information off-ledger.  G4 failure removes the finite tested resolution of the comparison.  G5 failure makes the route inverse or target-fitted.  G6 failure leaves mismatch unaccounted.  In each case the internal trace closure of $T$ is untouched, while the physical export continuation is not proved.  \qedhere
\end{proof}

\begin{corollary}[Route completion is sufficient for export-candidate status]
If $T$ is locally trace-closed and a route tuple
\[
        \mathcal R=(T,S,\Phi,\Lambda,\Sigma,\mathcal A,\mathcal C_{\rm res})
\]
satisfies G1--G6 with finite cost and without target reuse, then $T$ may be promoted to $[P_{\rm export\ candidate}]$ for $O_S$.
\end{corollary}

\begin{proof}
The completed route tuple supplies the typed codomain, the continuation witness, the on-ledger route commitments, the finite tested uncertainty structure, the independence audit, and the residual ledger.  These are precisely the ingredients needed for the APF comparison continuation $T\to O_S$ to be meaningful at export-candidate level.  Publication review or further physical calibration may still be required for $[P_{\rm physical\ final}]$, but the APF export-candidate predicate is satisfied.  \qedhere
\end{proof}

\begin{center}
\renewcommand{\arraystretch}{1.25}
\begin{tabular}{p{0.16\linewidth}p{0.36\linewidth}p{0.38\linewidth}}
\hline
\textbf{Gate} & \textbf{APF source} & \textbf{Derived necessity}\\
\hline
G1 Codomain & Contextual continuation; meaning by continuation profile & $O$ must be typed as $O_S$ before $T\to O_S$ is well formed.\\
G2 Transport & Primitive continuation judgment & Export requires a witness $\Phi:T\mapsto O_S$ or identity admission.\\
G3 Ledger & Specific finite commitments; closed ledger & Constants, scales, conventions, and normalizations used by $\Phi$ must be on-ledger.\\
G4 Covariance & Finite cost-bearing tested distinctions & Comparison requires a finite resolution/covariance structure propagated through the route.\\
G5 No-smuggling & Source-independent continuation; closed ledger & Target reuse collapses prediction into inverse fitting.\\
G6 Residuals & Closed ledger; residual as distinction & Mismatch must be assigned to named route/scheme/covariance channels.\\
\hline
\end{tabular}
\end{center}

\paragraph{Bottom line.}
In APF language, a physical observable is not a trace scalar.  It is a trace anchor plus an admissible export continuation into a declared physical codomain.  G1--G6 are the six typed components of that continuation.
